\section{Rectangle Cutting}
\label{sec:cses-rectangle-cutting}

\problemstatement{Given an $a\times b$ rectangle, repeatedly cut it with
straight horizontal or vertical cuts (each cut splits a piece into two,
along integer coordinates) until only squares remain. Find the minimum
number of cuts needed.}

\begin{patternbox}
\emph{``Minimum cuts to reduce a rectangle to squares''} is an interval
DP in two dimensions. Let $\textit{dp}[h][w]$ be the fewest cuts for an
$h\times w$ rectangle; a first cut splits it into two smaller rectangles
whose costs are already known.
\end{patternbox}

\subsection*{State and transition}

$\textit{dp}[h][w]$ = minimum cuts for an $h\times w$ piece. If $h=w$ it
is already a square: $\textit{dp}[h][h]=0$. Otherwise the first cut is
either horizontal (splitting height into $i$ and $h-i$) or vertical
(splitting width into $j$ and $w-j$), each costing one plus the two
sub-costs:
\[
  \textit{dp}[h][w]=1+\min\!\Big(
  \min_{1\le i<h}\!\big(\textit{dp}[i][w]+\textit{dp}[h-i][w]\big),\
  \min_{1\le j<w}\!\big(\textit{dp}[h][j]+\textit{dp}[h][w-j]\big)\Big).
\]

\begin{ideabox}[Try Every First Cut]
As in all interval DP, we do not guess the optimal cut -- we try every
possible first cut and take the minimum, trusting that sub-rectangles are
solved optimally. Squares terminate the recursion at zero cost.
\end{ideabox}

\subsection*{Algorithm}

\begin{enumerate}
  \item For all $s$, $\textit{dp}[s][s]=0$.
  \item For increasing $h,w$ with $h\ne w$: take the best over all
        horizontal splits $i$ and vertical splits $j$, add $1$.
  \item Answer $\textit{dp}[a][b]$.
\end{enumerate}

\subsection*{Dry run}

\dryrun{$a=3$, $b=5$.}

\begin{itemize}
  \item $\textit{dp}[3][3]=0$. $\textit{dp}[2][3]=1+\min$ over cuts
        $=1+\textit{dp}[2][2]+\textit{dp}[2][1]$-type options $=2$ (cut
        off a $2\times2$, then split the $2\times1$).
  \item Working up, $\textit{dp}[3][5]=3$ -- e.g.\ cut a $3\times5$ into
        $3\times3$ and $3\times2$, then the $3\times2$ into a $2\times2$
        and pieces, totalling three cuts.
\end{itemize}

\subsection*{C++ implementation}

\begin{lstlisting}
#include <bits/stdc++.h>
using namespace std;

int main() {
    int a, b;
    cin >> a >> b;
    vector<vector<int>> dp(a + 1, vector<int>(b + 1, 0));

    for (int h = 1; h <= a; h++)
        for (int w = 1; w <= b; w++) {
            if (h == w) { dp[h][w] = 0; continue; }
            int best = INT_MAX;
            for (int i = 1; i < h; i++)                 // horizontal cut
                best = min(best, dp[i][w] + dp[h - i][w]);
            for (int j = 1; j < w; j++)                 // vertical cut
                best = min(best, dp[h][j] + dp[h][w - j]);
            dp[h][w] = best + 1;
        }
    cout << dp[a][b] << "\n";
    return 0;
}
\end{lstlisting}

\begin{complexitybox}
\textbf{Time Complexity.} $O(ab\,(a+b))$ -- each of the $ab$ states tries
$O(a+b)$ cut positions. \textbf{Space Complexity.} $O(ab)$.
\end{complexitybox}

\begin{takeawaybox}
Interval DP in 2D: enumerate every first cut, recurse on the two pieces,
take the minimum. The ``try all split points'' pattern recurs in Removal
Game and any problem where a structure is optimally divided into
independent parts.
\end{takeawaybox}
